% HD-N15. GPT-6 Astra Pro, inherited effort, /root/imprimitive.
% 21 November 2026. Explicit character maps and effective degree bounds.
\subsection{HD-N15: explicit effective suppliers on Fermat varieties of degree twelve}
\label{sec:HD-N15}

Write $F_m^n\subset\mathbf P^{n+1}$ for
$\sum_{j=0}^{n+1}x_j^m=0$, and let $H$ be its hyperplane class.
The diagonal group is $G_m^n=\mu_m^{n+2}/\mu_m$.
A primitive character is a tuple
$\alpha=(a_0,\ldots,a_{n+1})$ with every $a_j\ne0$ in
$\mathbf Z/m$ and $\sum a_j=0$.
Its eigenspace $V(\alpha)$ in primitive middle cohomology has
dimension one.  For $n=2p$, its rational Galois orbit is of
Hodge type $(p,p)$ precisely when
\[
 \sum_j\frac{\langle ta_j\rangle}{m}=p+1
 \quad\text{for every }t\in(\mathbf Z/m)^\times,
 \qquad 1\leq\langle ta_j\rangle<m.
                                                        \tag{N15.1}
\]
These are Shioda's Theorem I, p.176, in
\emph{The Hodge conjecture for Fermat varieties},
Math. Ann.\ \textbf{245} (1979), 175--184,
\href{https://doi.org/10.1007/BF01428804}{doi:10.1007/BF01428804}.

The following bounds concern actual effective integral cycles.
Set
\begin{equation}
 \begin{split}
 M&=16(4\cdot12^3)^2=764411904,\\
 B&=12(M+10)=9172942968,\\
 E_{2p}&=B^p2^{p^2}\quad(p\geq1),\qquad E_0=1.
 \end{split}                                           \tag{N15.2}
\end{equation}
Every Hodge character of $F_{12}^{2p}$ has nonzero projection
from an effective irreducible $p$-dimensional cycle of
hyperplane degree at most $E_{2p}$.

More generally, let
\[
 V=\prod_{j=1}^R F_{12}^{r_j},\qquad r_j\geq1,
 \qquad S=\sum_jr_j\text{ even},\qquad
 N=S+2(R-1),\quad k=N+S.
\]
For each Hodge character in
$\bigotimes_jH^{r_j}_{\rm prim}(F_{12}^{r_j})$,
there is an effective irreducible cycle of dimension $S/2$
with nonzero projection to that character and degree, for
$H_V=\sum_jH_j$, at most
\begin{equation}
 \begin{split}
 C_{12}(S,R)={}&2(k+1)(k!12^{R+1})^{k+1}
       \bigl(2^{S+R-1}S!12^R\bigr)
       \bigl(k!12^R E_N\bigr)\\
 ={}&2^{S+R+(N/2)^2}(k+1)(k!)^{k+2}S!\,
       12^{(R+1)(k+1)+2R}B^{N/2}.
 \end{split}                                           \tag{N15.3}
\end{equation}
All the cycles can be chosen to meet the locus where every
homogeneous coordinate is nonzero.  For factors of degrees
$m_j\mid12$, $m_j\geq2$, the corresponding bound is
$12^{S/2}C_{12}(S,R)$, by the explicit coordinate-power morphisms
at the end of the proof.  Zero-dimensional factors are supplied
by points and may be removed before applying the formula.

\paragraph{1. A uniform effective curve on the Fermat surface.}
Looijenga, \emph{Fermat varieties and the periods of some
hypersurfaces}, Adv. Stud. Pure Math.\ \textbf{58} (2010),
47--67, Corollary 2.2 and the following pairing formula,
gives an integral generator $e$ of the primitive homology
as a $\mathbf Z[G_{12}^2]$-module.  The
\href{https://webspace.science.uu.nl/~looij101/fermatperiod4.pdf}
{author version}, pp.5--7, proves this from the affine
vanishing lattice and its integral quotient.  Under integral
Poincar\'e duality use the same letter for the corresponding
primitive class in $H^2(F_{12}^2,\mathbf Z)$.
The group-ring-valued intersection form is, up to sign,
\[
                         \prod_{j=0}^3(1-u_j).
\]
In particular every pairing between two orbit generators is
an integer of absolute value at most $16$: expansion has
sixteen terms, even when some terms represent the same group
element.  This estimate does not require a choice of an
orthogonal integral basis.

Fix a Galois Hodge orbit $\mathcal O$ of primitive characters.
It has at most $\varphi(12)=4$ elements.  The integral operator
\[
 P_{\mathcal O}
  =\sum_{g\in G_{12}^2}
       \left(\sum_{\chi\in\mathcal O}\chi(g)^{-1}\right)g
                                                        \tag{N15.4}
\]
has integer coefficients: the inner sum is both an algebraic
integer and invariant under the cyclotomic Galois group.
Their absolute values are at most four, and their total
absolute sum is at most $4\cdot12^3$.
Character orthogonality shows that $P_{\mathcal O}$ acts by
$12^3$ on each eigenspace in $\mathcal O$ and by zero on
the other eigenspaces.  The generator $e$ has nonzero projection
to every primitive eigenspace: otherwise its entire group-ring
span would have zero projection there.  Thus
\[
                D=P_{\mathcal O}e
\]
is a nonzero integral primitive class of type $(1,1)$, and
\[
                         D^2\geq-M.                    \tag{N15.5}
\]
Indeed the absolute value of the expanded self-pairing is
at most sixteen times the square of the coefficient mass
in (N15.4).

The integral Lefschetz $(1,1)$ theorem supplies a line bundle
with first Chern class $D$.  Put $a=M+10$ and $L=D+aH$.
Adjunction gives $K_{F_{12}^2}=8H$, and $H^2=12$.
The hypersurface cohomology sequence gives
$h^1(\mathcal O)=0$ and
$h^2(\mathcal O)=h^0(\mathcal O(8))=\binom{11}{3}=165$,
so $\chi(\mathcal O)=166$.
Since $D\cdot H=0$,
\[
 \chi(L)=166+\frac12\bigl(D^2+12a^2-96a\bigr)>0.
\]
The strict positivity follows already by replacing $D^2$ by
$-M$ and $a$ by $M+10$.
Moreover $(K-L)\cdot H=12(8-a)<0$, so $K-L$ has no
nonzero section.  Serre duality and Riemann--Roch therefore
give $h^0(L)\geq\chi(L)>0$.
The resulting effective divisor has degree $12a=B$ and the
same nonzero $\mathcal O$-projection as $D$; the hyperplane
class has trivial character.  At least one irreducible
component has nonzero projection to a specified character in
$\mathcal O$, and its degree is at most $B$.
Since its class is rational, nonvanishing holds on the whole
Galois orbit.  This establishes the surface input, with a
stronger bound than $E_2=2B$.

For $F_{12}^0$, any one of the twelve points has nonzero
projection to every nontrivial character of the regular
permutation representation; its degree is one.

\paragraph{2. The two exact character operations.}
The following specific published combinatorial input is used:
Shioda verifies $(P_{12}^n)$ for every even $n$ on p.182.
Together with its definitions on p.180, this says that every
Hodge character in even dimension $n>2$, after a permutation
of coordinates, has one of the following two decompositions:
\begin{align*}
 \alpha&=\beta*\gamma,
     &r+s&=n-2,&r,s&\geq0\text{ even};\\
 \alpha&=\beta\#\gamma,
     &r+s&=n,&r,s&>0\text{ even}.
\end{align*}
Here $\beta$ and $\gamma$ are Hodge characters of
$F_{12}^r$ and $F_{12}^s$.  The operation $*$ concatenates
all their coordinates.  For $\#$, write
$\beta=(b_0,\ldots,b_{r+1})$,
$\gamma=(c_0,\ldots,c_{s+1})$, require
$b_{r+1}+c_{s+1}=0$, and delete these last two coordinates
before concatenation.  They are nonzero coordinates, as required
for primitive characters.
This is the unconditional degree-twelve instance of Shioda's
verified combinatorics; no conjectural assertion for other
degrees is used.

We describe the actual geometric maps, so that the degree
induction also verifies survival of the character.
For the concatenation operation, put
\[
 E=\mathbf P_{F_{12}^r\times F_{12}^s}
       (\mathcal O(-H_1)\oplus\mathcal O(-H_2)),
 \qquad p:E\longrightarrow F_{12}^r\times F_{12}^s.
\]
We use the projectivization of lines.  There is a morphism
\[
 q:E\longrightarrow F_{12}^{r+s+2},
 \qquad [\lambda x,\mu y]\longmapsto
                   [\lambda x_0:\cdots:\lambda x_{r+1}:
                    \mu y_0:\cdots:\mu y_{s+1}].
                                                        \tag{N15.6}
\]
The Fermat equation vanishes separately on the two blocks.
The character homomorphism is restriction to the two coordinate
blocks, and its pullback sends $(\beta,\gamma)$ to
$\beta*\gamma$.
Shioda's Theorem II, pp.178--179, formulas (2.4), (2.7)
and (2.14), proves that $q_*p^*$ is nonzero on this
one-dimensional tensor character space; his map on that
summand is exactly $12q_*p^*$.

Choose effective irreducible middle-dimensional cycles
$Z_r,Z_s$ with nonzero projections to $\beta,\gamma$.
Then $q_*p^*[Z_r\times Z_s]$ is the effective join cycle.
Its character projection is nonzero, because the source
character has a nonzero tensor coefficient and the displayed
map is nonzero on its one-dimensional eigenspace.  No other
source character has that target character under the injective
pullback of character groups, so there is no cancellation with
the other components of the cycle class.
The degree of this join is
\[
                     \deg Z_r\,\deg Z_s.             \tag{N15.7}
\]
For a direct degree proof, write $\xi=q^*H$ and use the
projective-bundle formula
$p_*\xi^{r/2+s/2+1}=\sum_{u+v=r/2+s/2}H_1^uH_2^v$.
Only the term $(u,v)=(r/2,s/2)$ integrates nontrivially
on $Z_r\times Z_s$.

For the operation $\#$, choose $\epsilon^{12}=-1$ and use
the rational map
\begin{equation}
 \phi(x,y)=
 [x_0y_{s+1}:\cdots:x_ry_{s+1}:
       \epsilon x_{r+1}y_0:\cdots:\epsilon x_{r+1}y_s]
 :F_{12}^r\times F_{12}^s\dashrightarrow F_{12}^{r+s}.
                                                        \tag{N15.8}
\end{equation}
Its base locus is the smooth codimension-two product
$Y=\{x_{r+1}=y_{s+1}=0\}=F_{12}^{r-1}\times F_{12}^{s-1}$.
Blowing it up gives
$\pi:\widetilde X\to F_{12}^r\times F_{12}^s$ and a
morphism $\psi:\widetilde X\to F_{12}^{r+s}$.
Shioda's diagram (2.1), p.177, identifies this morphism as
a degree-twelve quotient followed by a blowdown and proves
that the restriction of $\psi_*\pi^*$ is a nonzero
isomorphism on the $(\beta,\gamma)$ character summand.
The equivariance map sends the target character precisely
to $(\beta,\gamma)$: the exponent of $x_{r+1}$ is the sum
of the surviving $\gamma$ coordinates, namely
$-c_{s+1}=b_{r+1}$, and the other last exponent is checked
in the same way.

There is an explicit check for using strict effective transforms
in place of the cohomological pullback.  Every irreducible
cycle with nonzero primitive-character projection meets the
all-coordinate-open locus.  Indeed, a cycle contained in
one coordinate hyperplane is fixed pointwise by the corresponding
coordinate $\mu_{12}$, whereas the character is nontrivial
on that group.  An irreducible cycle contained in the union
of these hyperplanes is contained in one of them.
Thus $Z=Z_r\times Z_s$ is not contained in $Y$, and its
strict transform $\widetilde Z$ is effective of the same
dimension.  The blowup cohomology formula gives
\[
 [\widetilde Z]-\pi^*[Z]
       \in j_*p_Y^*H^{r+s-2}(Y).
\]
Scaling $x_{r+1}$ fixes $Y$ pointwise and acts trivially
on this entire summand; it also preserves the tautological
class of the exceptional projective bundle.  But the source
character has nonzero exponent $b_{r+1}$ under this scaling.
Hence the difference has zero $(\beta,\gamma)$ projection.
It follows that the effective cycle $\psi_*[\widetilde Z]$
has the same nonzero target-character projection as
$\psi_*\pi^*[Z]$.

The map is finite over the all-coordinate-open locus, so the
pushforward has the required dimension.  Its hyperplane degree
is bounded by
\begin{equation}
 \deg\psi_*[\widetilde Z]
 \leq \binom{(r+s)/2}{r/2}\deg Z_r\deg Z_s
 \leq2^{(r+s)/2}\deg Z_r\deg Z_s.                 \tag{N15.9}
\end{equation}
To justify the first inequality, the resolved linear system has
$\psi^*H=\pi^*(H_1+H_2)-E_{\rm exc}$.
Both $\psi^*H$ and $\pi^*(H_1+H_2)$ are nef.  On the strict
transform their difference is an effective Cartier divisor.
The difference of their top powers is a sum of its intersections
with products of nef divisors, all nonnegative.  Projection
to $Z_r\times Z_s$ gives the binomial coefficient in (N15.9).
Passing to an irreducible component with nonzero character
projection only decreases the degree.  Coordinate permutations
preserve all degrees and give the required permuted character.

\paragraph{3. The universal induction bound.}
The base cases were proved in section 1.
For a join in dimension $2p$, let $r=2a,s=2b$ with
$a+b=p-1$.  Then
\[
 E_{2a}E_{2b}=B^{p-1}2^{a^2+b^2}
                         \leq B^p2^{p^2}=E_{2p}.
\]
For a $\#$ operation, $a,b\geq1$ and $a+b=p$.
The degree bound is
\[
 2^pE_{2a}E_{2b}=B^p2^{a^2+b^2+a+b}
                         \leq B^p2^{(a+b)^2}=E_{2p},
\]
because $2ab\geq a+b$ for positive integers $a,b$.
All dimensions used in these two operations are smaller than
$2p$.  This proves the bound (N15.2) for every even dimension.
The nonzero-character argument at each step also proves the
all-coordinate-open assertion.

\paragraph{4. The product incidence and its exact character adjoint.}
Let $V,S,R,N$ be as in the statement and write
\[
 \mathcal E=\bigoplus_{j=1}^R\mathcal O_V(-H_j),\qquad
 E=\mathbf P_V(\mathcal E),\qquad p:E\to V.
\]
Concatenating the $R$ blocks gives a morphism
$q:E\to F_{12}^N$, by the same formula as (N15.6).
The map
\[
                (q,p):E\longrightarrow F_{12}^N\times V
                                                        \tag{N15.10}
\]
is a closed immersion: it is the projective subbundle associated
to the direct sum of the $R$ tautological coordinate lines in
the trivial coordinate bundle on $V$.
Denote its image by $\Gamma$.  In particular it is a closed
algebraic correspondence of dimension $S+R-1$.

The concatenated character
$\alpha=\beta_1*\cdots*\beta_R$ has length $N+2$.
For every Galois conjugate, the product Hodge condition is
\[
 \sum_j|t\beta_j|=S/2+R=N/2+1.
\]
Thus $\alpha$ is exactly a Hodge character of $F_{12}^N$.
The map
\[
 T=q_*p^*:H^S(V)\longrightarrow H^N(F_{12}^N)
                                                        \tag{N15.11}
\]
is nonzero on the tensor character space in question.
Here is an actual comparison with the iterated binary joins
of section 2.  Iterating their projective bundles gives a
proper space mapping to $E$ by multiplying the successive
scalar coefficients in each block.  Over the locus where all
block coefficients are nonzero, this is an isomorphism; it is
therefore proper birational.  Its two maps to $V$ and
$F_{12}^N$ factor through $p$ and $q$.
At each composition the next projective-bundle projection is
flat, so proper base change identifies composition of the
cohomological correspondences with the map of this iterated
bundle.  Proper pushforward along the birational map, followed
by the projection formula, gives exactly (N15.11).
Consequently (N15.11) is the composition of the nonzero binary
character maps already identified.  Each intermediate character
space is one-dimensional; hence this composition is nonzero.

Its Poincar\'e adjoint is the actual map
\[
 T^\dagger=p_*q^*:H^N(F_{12}^N)\longrightarrow H^S(V).
                                                        \tag{N15.12}
\]
The pairing between the $\alpha$ and $\bar\alpha$ lines is
perfect, as is the product pairing on the two inverse tensor
characters.  Applying the adjunction identity
\[
 \int_{F_{12}^N}T(x)\smile y
                   =\int_Vx\smile T^\dagger(y)
\]
to the nonzero map on the inverse character shows that
$T^\dagger$ is nonzero on the $\alpha$ line.
There is no contribution to this tensor character from any
other target character, by equivariance for restriction to
the $R$ coordinate blocks.

Choose an effective irreducible middle cycle $Y_N$ on
$F_{12}^N$ with nonzero $\alpha$ projection and degree at
most $E_N$.  Then $T^\dagger[Y_N]$ has the desired nonzero
tensor-character projection.  We now realize this class by
bounded effective cycles, using intersection in a smooth
ambient product.  No effectiveness assertion about a refined
pullback is needed.

\paragraph{5. Quantitative intersection and an effective component.}
Embed
\[
 X=F_{12}^N\times V
\]
by the product hyperplane line bundle
$H_X=H_N+\sum_jH_j$.  Its dimension is $k=N+S$ and its
degree satisfies
\begin{equation}
 \Delta_X=\frac{k!}{N!\prod_jr_j!}\,12^{R+1}
                             \leq k!12^{R+1}.          \tag{N15.13}
\end{equation}
Put $t=S+R-1$ and $\xi=q^*H_N$ on $E$.
The projective-bundle formula says
$p_*\xi^{R-1+j}=h_j(H_1,\ldots,H_R)$, the complete
homogeneous polynomial.  Its coefficients are at most those
of $(\sum H_i)^j$, and all the divisor classes are nef.
Consequently
\begin{align*}
 \deg_{H_X}\Gamma
 &=\int_E(\xi+p^*H_V)^t\\
 &\leq\sum_{j=0}^S\binom{t}{R-1+j}\int_VH_V^S\\
 &\leq2^{S+R-1}\frac{S!}{\prod_jr_j!}\,12^R
 \leq2^{S+R-1}S!12^R.                              \tag{N15.14}
\end{align*}
Likewise
\begin{equation}
 \deg_{H_X}(Y_N\times V)
 =\frac{(N/2+S)!}{(N/2)!\prod_jr_j!}\,
                       12^R\deg Y_N
 \leq k!12^R E_N.                                  \tag{N15.15}
\end{equation}

The quantitative intersection result used here is Truong,
\emph{Relative dynamical degrees of correspondences over a
field of arbitrary characteristic}, J. Reine Angew. Math.\
\textbf{758} (2020), 139--182,
\href{https://doi.org/10.1515/crelle-2017-0052}
{doi:10.1515/crelle-2017-0052}.
The precise checked locator is Lemma 2.2 and its full proof
on p.13 of the
\href{https://arxiv.org/pdf/1605.05049}{author version}.
For two irreducible cycles $A,B$ on a smooth projective
$k$-fold of degree $\Delta$, their Chow intersection is a
difference $A_+-A_-$ of effective cycles.  The proof gives
the constant $k\Delta^k$; we use the larger
$(k+1)\Delta^{k+1}$.
Indeed Roberts' moving construction in that proof has at most
$k$ successive cones and one final residual cycle.  Each
cone has the same degree as its base, and each intersection
with the ambient $X$ multiplies the available bound by at
most $\Delta$.  Intersecting the properly moved terms with
the other cycle gives at most $k+1$ effective terms of degree
at most $\Delta^k\deg A\deg B$.
This proves the conservative bound we use directly from
the cited construction.

Apply it to $\Gamma$ and $Y_N\times V$ in this smooth $X$.
Their intersection has dimension $S/2$, since
\[
 (S+R-1)+(N/2+S)-(N+S)=S/2.
\]
There are effective cycles $A_+,A_-$ of that dimension with
\[
 [\Gamma]\cdot[Y_N\times V]=[A_+]-[A_-]
\]
in the Chow group, and each of their degrees is at most
\[
 (k+1)(k!12^{R+1})^{k+1}
       (2^{S+R-1}S!12^R)(k!12^RE_N).
\]
Pushing them to $V$ gives effective cycles, discarding any
component on which dimension drops.  This pushforward does
not increase degree, because the pullback of $H_V$ and its
complement $H_N$ in $H_X$ are nef.  Its difference represents
exactly (N15.12) applied to $[Y_N]$, by the definition of
the graph correspondence and the projection formula.
That class has a nonzero tensor-character projection.
At least one irreducible component of one of the two effective
pushforwards has nonzero projection.  Its degree is bounded
by the sum of the two bounds, which is (N15.3).
The coordinate-group argument from section 2 shows that this
component meets the all-coordinate-open locus.

\paragraph{6. Smaller degrees, points and the rational orbit.}
For $m\mid12$, put $d=12/m$.  The actual finite morphism
\[
 f_m:F_{12}^r\longrightarrow F_m^r,
                 \qquad[x_0:\cdots:x_{r+1}]
                        \longmapsto[x_0^d:\cdots:x_{r+1}^d]
                                                        \tag{N15.16}
\]
has $f_m^*H_m=dH_{12}$.  It pulls a primitive character
$(a_j)\bmod m$ to $(da_j)\bmod12$, still with every
coordinate nonzero.  Its pullback on that one-dimensional
cohomology space is nonzero, because
$f_{m*}f_m^*=\deg(f_m)$; its pushforward is nonzero there
for the same reason.  The Hodge criterion (N15.1) is
preserved by these character formulas.  Products of these
maps therefore carry the bounded degree-twelve supplier
to an effective detecting cycle for the prescribed smaller
degrees.  For a cycle of dimension $S/2$ its degree is at most
\[
 \int_Z\left(\sum_jd_jH_j\right)^{S/2}
 \leq12^{S/2}\int_Z\left(\sum_jH_j\right)^{S/2}.
\]
This proves the final degree factor in the statement.

A factor $F_m^0$ is a set of $m$ points.  A single point
has nonzero projection to each nontrivial character of its
regular permutation representation.  Taking its external
product with a supplier on the positive-dimensional factors
preserves the required character and the degree.  When every
factor is zero-dimensional, a point in the product has
degree one and supplies all the tensor characters in question.

Finally, for a rational Galois character orbit, a rational
cycle nonzero at one character is nonzero at each conjugate.
Its orbit under the finite diagonal group has projected span
equal to the full rational irreducible character module:
after extension to the cyclotomic field, the character
projectors extract each of its nonzero one-dimensional
components.  Every translate is an effective cycle of the
same degree.  Thus the bounds supply all the relevant
character directions, not only one unspecified vector.

\paragraph{Source and check scope.}
The integral generator and pairing were checked in Looijenga's
original author text, including the proof of Corollary 2.2.
Shioda's printed pp.176--182 were read, with visual checks of
the diagram, character maps and formulas on pp.177--179 and
the degree-twelve combinatorial verification on p.182.
The complete proof of Truong's Lemma 2.2 was independently
read at its stated author-version locator, and its publication
was checked against the publisher record.
The geometric degree bounds and character-survival argument
above are derived here from those explicit inputs.  No
unrefereed Fermat supplier theorem is used.
The accompanying exact arithmetic certificate checks the
surface constants and finite character projectors; the
all-dimensional induction and the actual correspondence
maps are proved in the text rather than inferred from samples.
